\scanpage{74}
Of course, $Q(A)$ is a $\handwritten{\lambda}$-ring, and its $\handwritten{\lambda}$-structure can be read off from the ring structure, since we will have
\[
\handwritten{\lambda}_t(\handwritten{\varphi}_a)=1+\handwritten{\varphi}_a t\ ,
\]
On the other hand, a homomorphism of augmented rings $Q(A)\handwritten{\to}\underline{Z}\times(1+\overset{\handwritten{\scriptscriptstyle\wedge}}{A}^{\handwritten{+}})$ is automatically a $\handwritten{\lambda}$-homomorphism. I do not think there is any reason to say more about $\handwritten{\lambda}$-structures (except perhaps in a remark that the filtration put on $Q(A)$ is the one deduced from its $\handwritten{\lambda}$-structure).

2) \uline{Case where} $A$ \uline{is a local ring}, with residue field $k$ such that $\mathrm{car}\ k\handwritten{\ne}2$. It does seem that all that precedes works in this case; of course, $H^{\handwritten{\bullet}}(A)$ denotes the cohomology of the ``fundamental group'' of $\mathrm{Spec}(A)$, in the sense of my seminar. One would then have to write everything up in the setting: quadratic forms over a local ring. N.B. such a form is said to be nondegenerate if the reduction is nondegenerate; in principle, it \deleted{\unsure{would have to be ca|would have to be that}} should be possible in a fairly trivial way to lift the usual constructions over $k$, to constructions over $A$.

Of course, it would be very pleasant to go beyond the minimum programme and arrive at a theory of Stiefel-Whitney classes for quadratic Modules on an arbitrary \inserted{pre}scheme $S$, \deleted{\unsure{One could commence|One could commen}} assuming if need be that the residue characteristics are $\handwritten{\ne}2$. \deleted{trit} As a test, one can always try to prove the invariance of the trivial definition relative to a \inserted{quadratic} Module which splits into \inserted{quadratic} modules of rank 1.

As a bonus for the reader, one can in any case give the compatibility of the purely algebraic definition of Stiefel-Whitney classes, and of the topological definition, when $A$ is the local ring\inserted{ of a point $s$} of a scheme of finite type\inserted{$S$} over $\underline{C}$. A quadratic module over $A$ defines a quadratic module on a neighborhood of $S$, hence $S$ if one wishes, i.e. an algebraic principal homogeneous bundle under the orthogonal group scheme. By Gaga there corresponds to it an analytic principal bundle, hence topological, whence topological Stiefel-Whitney classes in $H^{\handwritten{\bullet}}(S_{\mathrm{top}},\underline{Z}/2\underline{Z})$. Now one has a canonical homomorphism,which you will define:
\[
H^{\handwritten{\bullet}}(S_{\mathrm{top}},\underline{Z}/2\underline{Z})\handwritten{\longrightarrow}H^{\handwritten{\bullet}}(A,\underline{Z}/2\underline{Z})\ ,
\]
and one sees at once that the images of the topological characteristic classes are the algebraic characteristic classes, (for it suffices to verify this in rank 1).
% archive sheet number: 73
